%%
%% fall-lecture04.tex
%% 
%% Made by Alex Nelson <pqnelson@gmail.com>
%% Login   <alex@lisp>
%% 
%% Started on  2026-10-13T13:05:04-0700
%% Last update 2026-10-13T13:05:04-0700
%% 

\lecture{}

\begin{corollary}[{\cite[Th~3.3]{milnor1974characteristic}}]
Let $\xi$ be a sub-bundle of $\eta$. Then $\xi\oplus\xi^{\perp}\iso\eta$.
\end{corollary}

\begin{example}
Let $M$ be a smooth manifold. Suppose $M\into K$ is an embedding,
where $K$ is a Riemannian manifold. Then $\T M$ is a sub-bundle of $\T K|_{M}$
(if $i\colon M\into K$ is the embedding, $i^{*}\T K =: TK|_{M}$). Then
$(\T M)^{\perp}=NM$ is the normal bundle of $M$.
\end{example}

\begin{example}
For $K=\sphere{2}$ and $M=\sphere{1}$ is the equator of $K$, then
$NM\iso\sphere{1}\iso\RR$ is trivial because we have a nowhere
vanishing section (take the tangent vector pointing Northwards, its
span is precisely $\RR$). Observe $\T\sphere{2}|_{\sphere{1}}=\T\sphere{1}\oplus(\sphere{1}\times\RR)$.
\end{example}

\begin{node}
Let $V$, $W$ be finite-dimensional [real] vector spaces. We have
constructions like the vector space of linear transformations $\hom(V,W)$,
the tensor product $V\otimes W$,
the space of linear functionals $V^{*}$,
and the exterior power $\extp^{k}V$.
These constructions may be generalized to vector
bundles. We get $\hom(\xi,\eta)$, $\xi\otimes\eta$, $\extp^{k}\xi$.
\end{node}

\begin{notation}
We denote by $\cat{C}$ the category
\begin{itemize}
\item \textsc{Objects}: all finite-dimensional vector spaces over $\RR$
\item \textsc{Morphisms}: all $\RR$-linear isomorphisms between vector spaces
\end{itemize}
\end{notation}

\begin{definition}
A \define{Continuous Covariant Functor} $T\colon\cat{C}\times\cat{C}\to\cat{C}$
is an assignment
\begin{enumerate}
\item for any objects $V,W\in\cat{C}$ we have $T(V,W)\in\cat{C}$ is a
  finite-dimensional [real] vector space
\item for any morphisms $f\colon V\to V'$ and $g\colon W\to W'$ in $\cat{C}$
we have a morphism $T(f,g)\colon T(V,W)\to T(V',W')$
\end{enumerate}
such that
\begin{enumerate}
\item $T(\id_{V},\id_{W})=\id_{T(V,W)}$
\item $T(f_{1}\circ f_{2},g_{1}\circ g_{2})=T(f_{1},g_{1})\circ T(f_{2},g_{2})$
\item $T(f,g)$ depends continuously on $f$ and $g$. (This makes sense
  since the set of all isomorphisms between two real vector spaces has
  a natural topology on it.)
\end{enumerate}
We can extend this definition to $T\colon\cat{C}\times(\cdots)\times\cat{C}\to\cat{C}$
any $k$-ary such assignment.
\end{definition}

\begin{theorem}
Let $\xi_{1}$, \dots, $\xi_{k}$ be real vector bundles over the same
base space $B$. Let $T\colon\cat{C}\times(\cdots)\times\cat{C}\to\cat{C}$
be a $k$-ary continuous covariant functor. Then we can define a new
bundle with fiber $F_{b}=T(F_{b}(\xi_{1}),\dots,F_{b}(\xi_{k}))$.
\end{theorem}

\begin{xca}
Prove this is a vector bundle.
\end{xca}

\begin{xca}
What constructions can be defined using continuous covariant functors?
\end{xca}

\subsection{Grassmannians and universal bundles}

\begin{node}
Let $G_{n}(\RR^{\infty})$ consist of all $n$-planes passing through
the origin in $\RR^{\infty}$. In particular, for $n=1$, this is just
$G_{1}(\RR^{\infty})=\RP^{\infty}$. This is going to be the universal
base space.
\end{node}

\begin{node}
Then $\gamma^{n}$ is an $\RR^{n}$-bundle over $G_{n}(\RR^{\infty})$.
We have
\begin{equation}
\gamma^{n}=\{(x,v)\in G_{n}(\RR^{\infty}\times\RR^{\infty}\mid x\in X\},
\end{equation}
and $\pi(x,v)=x$.
\end{node}

\begin{definition}
The \define{Grassmann Manifold} $G_{n}(\RR^{n+k})$ is the set of all
$n$-dimensional planes passing through the origin of $\RR^{n+k}$,
equipped with ``some topology'' (which we will now determine).
\end{definition}

\begin{definition}
An \define{$n$-Frame} in $\RR^{n+k}$ is an $n$-tuple that is linearly
independent $(v_{1},\dots,v_{n})\in(\RR^{n+k})^{n}$ and for all $i\neq j$
we have $v_{i}$ is linearly independeot of $v_{j}$.
\end{definition}

\begin{definition}
We define the \define{Stiefel Manifold}
$V_{n}(\RR^{n+k})$ to be the set of all $n$-frames in
$\RR^{n+k}$, considered as a subspace of $(\RR^{n+k})^{n}$ endowed
with the subspace topology from the standard topology on $(\RR^{n+k})^{n}$.
\end{definition}

\begin{fact}
We can obtain $q\colon V_{n}(\RR^{n+k})\to G_{n}(\RR^{n+k})$ which is
a quotient map sending $(v_{1},\dots,v_{n})\mapsto\Span(v_{1},\dots,v_{n})$.
So
\begin{equation}
G_{n}(\RR^{n+k})\iso V_{n}(\RR^{n+k})/\GL_{n}(\RR)
\end{equation}
[where we mod out by a change of coordinates]. Then we can equip
$G_{n}(\RR^{n+k})$ with the quotient topology:
$U\subset G_{n}(\RR^{n+k})$ is open if and only if $q^{-1}(U)$ is open.
\end{fact}

\begin{lemma}
$G_{n}(\RR^{n+k})$ is a compact topological manifold of dimension $nk$.
Moreover, the orthogonal complement operation
\begin{align*}
\perp\colon & G_{n}(\RR^{n+k})\to G_{k}(\RR^{n+k})\\
& Y\mapsto Y^{\perp}
\end{align*}
is a homeomorphism. (These are 2 claims.)
\end{lemma}